\documentclass[border=5pt,svgnames,tikz]{standalone}
\usepackage{fontspec}
\usepackage{ctex}
\setmainfont{Source Serif 4}
\setsansfont{Source Sans 3}
\setmonofont{Source Code Pro}
\setCJKmainfont[BoldFont=Source Han Sans CN]{Source Han Sans CN}
\usetikzlibrary{positioning, calc, arrows.meta}

\begin{document}
\begin{tikzpicture}[
    font=\sffamily,
    node distance=0pt,
    block/.style={draw, rectangle, minimum height=1.2cm, inner sep=0pt, outer sep=0pt, anchor=west},
    smallblock/.style={draw, rectangle, minimum height=1.2cm, minimum width=0.2cm, inner sep=0pt, outer sep=0pt, anchor=west},
    label text/.style={font=\sffamily\small, align=center}
]

% Row 3: Disk Structure
% MBR (Width 80 units -> 2.5cm)
\node[block, minimum width=2.5cm] (mbr) {MBR};

% Partition Table Entries (4 small blocks, Width 5 units -> 0.15cm)
\node[smallblock, right=0pt of mbr] (pt1) {};
\node[smallblock, right=0pt of pt1] (pt2) {};
\node[smallblock, right=0pt of pt2] (pt3) {};
\node[smallblock, right=0pt of pt3] (pt4) {};

% Partitions (4 large blocks, Width 150 units -> 4.5cm)
\node[block, minimum width=4.5cm, right=0pt of pt4] (p1) {};
\node[block, minimum width=4.5cm, right=0pt of p1] (p2) {};
\node[block, minimum width=4.5cm, right=0pt of p2] (p3) {};
\node[block, minimum width=4.5cm, right=0pt of p3] (p4) {};

% Row 2: Labels
\node[above=1cm of pt2, anchor=south, xshift=-0.5cm] (lbl_pt) {分区表};
\node[above=1cm of p2, anchor=south] (lbl_part) {磁盘分区};

% Arrows for Row 2
\draw[-{Stealth}] (lbl_pt) -- (pt2.north);
\draw[-{Stealth}] (lbl_part) -- (p1.north);
\draw[-{Stealth}] (lbl_part) -- (p2.north);
\draw[-{Stealth}] (lbl_part) -- (p3.north);
\draw[-{Stealth}] (lbl_part) -- (p4.north);

% Row 1: Entire Disk
\node[above=2.5cm of p2, anchor=center] (lbl_disk) {\Large 整个磁盘};
% Arrow spanning the whole disk
\draw[{Stealth}-{Stealth}, thick] ($(mbr.west)+(0, 2.5)$) -- ($(p4.east)+(0, 2.5)$);
% Vertical ticks for the arrow
\draw[thick] ($(mbr.west)+(0, 2.3)$) -- ($(mbr.west)+(0, 2.7)$);
\draw[thick] ($(p4.east)+(0, 2.3)$) -- ($(p4.east)+(0, 2.7)$);


% Row 4: Expanded Partition Details (Expanding Partition 2)
% Total width of expanded row: 80+80+125+125+80+125 = 615 units.
% Scale: 150 units = 4.5cm => 1 unit = 0.03cm
% 80 units = 2.4cm
% 125 units = 3.75cm
% Total width = 615 * 0.03 = 18.45cm
% Center of P2 is at x of P2.center.
% We want to center the expanded row roughly below P2, or match the SVG layout.
% In SVG, expanded row starts at x=26.5, P2 starts at 259.
% Let's just place it below and center it relative to the whole diagram or P2.
% Let's center it relative to P2 for better aesthetics, or span it across.

\node[block, minimum width=2.4cm, below=2.5cm of mbr.west, anchor=west, xshift=0.5cm] (boot) {引导块};
\node[block, minimum width=2.4cm, right=0pt of boot] (super) {超级块};
\node[block, minimum width=3.75cm, right=0pt of super] (free) {空闲空间管理};
\node[block, minimum width=3.75cm, right=0pt of free] (inode) {index 节点};
\node[block, minimum width=2.4cm, right=0pt of inode] (root) {根目录};
\node[block, minimum width=3.75cm, right=0pt of root] (files) {文件和目录};

% Expansion lines (From Partition 2)
\draw (p2.south west) -- (boot.north west);
\draw (p2.south east) -- (files.north east);

\end{tikzpicture}
\end{document}
